% GENERATED FILE. Edit canonical lessons, then run npm run generate:books.
% canonical-source-hash: fnv1a64:003892a9c47000d6
% canonical-lessons: GE-C141-zahlen, GE-C141-mieten, GE-C141-packen, GE-C141-reparieren, GE-C141-anfangen, GE-R141-first-pass-wanting-time-and-numbers, GE-R141-second-pass-verbs

\chapter{To pay, To rent, To pack, To repair, To start}
\label{ch:ge-to-pay-to-rent-to-pack-to-repair-to-start}
\hlchaptermodality{141}

\begin{chapteropening}
\textbf{By the end of this chapter:} \emph{I can say to pay, to rent, to pack, to repair and to start.}

Complete the last lesson of chapter 141: i can say to pay, to rent, to pack, to repair and to start.
\end{chapteropening}

\section[zahlen]{zahlen — to pay}

\label{lesson:GE-C141-zahlen}



\begin{quote}
Before the new one: say the German for to feel, then the German for to fall.
\end{quote}

\subsection*{You'll want to know: zahlen}
\textbf{zahlen} — \textquotedblleft{}to pay\textquotedblright{}.

\begin{grammarlens}[title={What you've built}]
One more word for this chapter.
\end{grammarlens}

\subsection*{Your turn}
\begin{itemize}
\raggedright
  \item \emph{Say it:} \emph{zahlen}
  \item \emph{Say it:} \emph{zahlen}, once more
  \item \emph{Say it:} say \emph{zahlen}
  \item \emph{Recall:} say the German for to feel, then the German for to fall, then say \emph{zahlen} again
\end{itemize}

\begin{tcolorbox}[breakable,colback=teal!4,colframe=teal!35!black,title={Before you move on}]
What does zahlen mean? (\textquotedblleft{}To pay\textquotedblright{}.) Say it once more.
\end{tcolorbox}
\section[mieten]{mieten — to rent}

\label{lesson:GE-C141-mieten}



\begin{quote}
Before the new one: say the German for to fall, then the German for to pay.
\end{quote}

\subsection*{You'll want to know: mieten}
\textbf{mieten} — \textquotedblleft{}to rent\textquotedblright{}.

\begin{grammarlens}[title={What you've built}]
One more word for this chapter.
\end{grammarlens}

\subsection*{Your turn}
\begin{itemize}
\raggedright
  \item \emph{Say it:} \emph{mieten}
  \item \emph{Say it:} \emph{mieten}, once more
  \item \emph{Say it:} say \emph{mieten}
  \item \emph{Recall:} say the German for to fall, then the German for to pay, then say \emph{mieten} again
\end{itemize}

\begin{tcolorbox}[breakable,colback=teal!4,colframe=teal!35!black,title={Before you move on}]
What does mieten mean? (\textquotedblleft{}To rent\textquotedblright{}.) Say it once more.
\end{tcolorbox}
\section[packen]{packen — to pack}

\label{lesson:GE-C141-packen}



\begin{quote}
Before the new one: say the German for to pay, then the German for to rent.
\end{quote}

\subsection*{You'll want to know: packen}
\textbf{packen} — \textquotedblleft{}to pack\textquotedblright{}.

\begin{grammarlens}[title={What you've built}]
One more word for this chapter.
\end{grammarlens}

\subsection*{Your turn}
\begin{itemize}
\raggedright
  \item \emph{Say it:} \emph{packen}
  \item \emph{Say it:} \emph{packen}, once more
  \item \emph{Say it:} say \emph{packen}
  \item \emph{Recall:} say the German for to pay, then the German for to rent, then say \emph{packen} again
\end{itemize}

\begin{tcolorbox}[breakable,colback=teal!4,colframe=teal!35!black,title={Before you move on}]
What does packen mean? (\textquotedblleft{}To pack\textquotedblright{}.) Say it once more.
\end{tcolorbox}
\section[reparieren]{reparieren — to repair}

\label{lesson:GE-C141-reparieren}



\begin{quote}
Before the new one: say the German for to rent, then the German for to pack.
\end{quote}

\subsection*{You'll want to know: reparieren}
\textbf{reparieren} — \textquotedblleft{}to repair\textquotedblright{}.

\begin{grammarlens}[title={What you've built}]
One more word for this chapter.
\end{grammarlens}

\subsection*{Your turn}
\begin{itemize}
\raggedright
  \item \emph{Say it:} \emph{reparieren}
  \item \emph{Say it:} \emph{reparieren}, once more
  \item \emph{Say it:} say \emph{reparieren}
  \item \emph{Recall:} say the German for to rent, then the German for to pack, then say \emph{reparieren} again
\end{itemize}

\begin{tcolorbox}[breakable,colback=teal!4,colframe=teal!35!black,title={Before you move on}]
What does reparieren mean? (\textquotedblleft{}To repair\textquotedblright{}.) Say it once more.
\end{tcolorbox}
\section[anfangen]{anfangen — to start}

\label{lesson:GE-C141-anfangen}



\begin{quote}
Before the new one: say the German for to pack, then the German for to repair.
\end{quote}

\subsection*{You'll want to know: anfangen}
\textbf{anfangen} — \textquotedblleft{}to start\textquotedblright{}.

\begin{grammarlens}[title={What you've built}]
One more word for this chapter.
\end{grammarlens}

\subsection*{Your turn}
\begin{itemize}
\raggedright
  \item \emph{Say it:} \emph{anfangen}
  \item \emph{Say it:} \emph{anfangen}, once more
  \item \emph{Say it:} say \emph{anfangen}
  \item \emph{Recall:} say the German for to pack, then the German for to repair, then say \emph{anfangen} again
\end{itemize}

\begin{tcolorbox}[breakable,colback=teal!4,colframe=teal!35!black,title={Before you move on}]
What does anfangen mean? (\textquotedblleft{}To start\textquotedblright{}.) Say it once more.
\end{tcolorbox}
\section[(dialogue)]{Wanting, time and numbers — a first pass}

\label{lesson:GE-R141-first-pass-wanting-time-and-numbers}



\begin{quote}
Say the words for to repair and to start.
\end{quote}

\subsection*{Your turn}
Say these aloud:

\begin{itemize}
\raggedright
  \item angry; bad, funny, serious, famous, interesting
  \item boring, difficult, finished, ready, safe; sure
  \item dangerous, wet, dry, cool, soft
  \item to want, a wish, to hope, to order, rather
\end{itemize}

\begin{tcolorbox}[breakable,colback=teal!4,colframe=teal!35!black,title={Before you move on}]
Angry; bad? (\textbf{böse}.) Funny? (\textbf{lustig}.) Serious? (\textbf{ernst}.) Famous? (\textbf{berühmt}.) Interesting? (\textbf{interessant}.) Boring? (\textbf{langweilig}.) Difficult? (\textbf{schwierig}.) Finished? (\textbf{fertig}.) Ready? (\textbf{bereit}.) Safe; sure? (\textbf{sicher}.) Dangerous? (\textbf{gefährlich}.) Wet? (\textbf{nass}.) Dry? (\textbf{trocken}.) Cool? (\textbf{kühl}.) Soft? (\textbf{weich}.) To want? (\textbf{wollen}.) A wish? (\textbf{der Wunsch}.) To hope? (\textbf{hoffen}.) To order? (\textbf{bestellen}.) Rather? (\textbf{lieber}.)
\end{tcolorbox}
\section[(dialogue)]{Verbs — a second pass}

\label{lesson:GE-R141-second-pass-verbs}



\begin{quote}
Say the words for the weekend and a date.
\end{quote}

\subsection*{Your turn}
Say these aloud:

\begin{itemize}
\raggedright
  \item the weekend, a date, a century, never, often
  \item sixty, seventy, eighty, ninety, a million
  \item to begin, to visit, to answer, to lay, to put
  \item to lie, to get dressed, to tell, to feel, to fall
  \item to pay, to rent, to pack, to repair, to start
\end{itemize}

\begin{tcolorbox}[breakable,colback=teal!4,colframe=teal!35!black,title={Before you move on}]
The weekend? (\textbf{das Wochenende}.) A date? (\textbf{das Datum}.) A century? (\textbf{das Jahrhundert}.) Never? (\textbf{nie}.) Often? (\textbf{oft}.) Sixty? (\textbf{sechzig}.) Seventy? (\textbf{siebzig}.) Eighty? (\textbf{achtzig}.) Ninety? (\textbf{neunzig}.) A million? (\textbf{die Million}.) To begin? (\textbf{beginnen}.) To visit? (\textbf{besuchen}.) To answer? (\textbf{antworten}.) To lay? (\textbf{legen}.) To put? (\textbf{stellen}.) To lie? (\textbf{liegen}.) To get dressed? (\textbf{sich anziehen}.) To tell? (\textbf{erzählen}.) To feel? (\textbf{fühlen}.) To fall? (\textbf{fallen}.) To pay? (\textbf{zahlen}.) To rent? (\textbf{mieten}.) To pack? (\textbf{packen}.) To repair? (\textbf{reparieren}.) To start? (\textbf{anfangen}.)
\end{tcolorbox}
